\documentclass[11pt]{article}
\usepackage[margin=1in]{geometry}
\usepackage{amsmath,amssymb,amsthm}
\usepackage{hyperref}
\usepackage[numbers,sort&compress]{natbib}

\newtheorem{theorem}{Theorem}
\newtheorem{lemma}{Lemma}
\newtheorem{corollary}{Corollary}
\newtheorem{definition}{Definition}

\title{Literature Review: Set Theoretic Estimation\\
\large Working notes toward a Lean mechanization}
\author{Project \texttt{thejeb}}
\date{\today}

\begin{document}
\maketitle

\begin{abstract}
Working notes reviewing set theoretic estimation (STE) from its unifying
statement \citep{combettes1993} through its topological-space,
cryptographic instantiation \citep{carlson2012}, oriented toward
identifying which results are ripe for machine-checked formalization in
Lean~4 / Mathlib. These notes are exploratory and cite the primary
sources; only claims discharged in the CI-checked \texttt{Ste} library
are asserted as proven.
\end{abstract}

\section{The set theoretic paradigm (Combettes 1993)}

Conventional estimation optimizes an objective function; the solvability
and realism of that optimization is a persistent difficulty
\citep{combettes1993}. STE replaces optimization with \emph{feasibility}.

Fix a \emph{solution space} $\Xi$ and a true object (estimandum) $h$.
Each piece of information $i\in I$ is a proposition on $\Xi$; at a chosen
confidence level it yields a \emph{property set}
\begin{equation}
  S_i = \{a\in\Xi \mid \Psi_i(a)\ge\psi_i\},
  \qquad \Psi_i:\Xi\to[0,1],
\end{equation}
the set of estimates consistent with information $i$ (Combettes Eq.~3).
The \emph{feasibility set} is their intersection,
\begin{equation}
  S = \bigcap_{i\in I} S_i \qquad\text{(Combettes Eq.~4).}
\end{equation}
Any $a\in S$ is a \emph{set theoretic estimate}. The formulation is
\emph{consistent} if $S\ne\emptyset$, \emph{fair} if $h\in S$, and
\emph{ideal} if $S=\{h\}$. When $\Xi$ is a real Hilbert space and each
$S_i$ is closed and convex, $S$ is closed convex and can be reached by
projection algorithms (POCS, \citealp{youla1982}; cyclic projections,
\citealp{bregman1967}; survey \citealp{bauschke1996}).

\paragraph{What we formalized.} The \emph{crisp} core --- where $\Psi_i$
is $\{0,1\}$-valued so $S_i$ is an ordinary subset --- is machine-checked
in \texttt{Ste.Basic}: feasibility set, membership, refinement of each
constraint, information monotonicity, the fair/consistent link, the
inconsistency$\Rightarrow$invalid-information self-diagnostic, and the
ideal-information identification theorem. Convexity of the feasibility
set is \texttt{Ste.Convex}. The graded/fuzzy case and the projection
\emph{algorithms} (as opposed to the geometry) remain open.

\section{Topological-space STE for decryption (Carlson 2012)}

\citet{carlson2012} makes two moves. First, a \emph{spatial} move: prior
STE works in a Hilbert / metric space and simplifies processing with
Optimal Bounding Ellipsoid (OBE) algorithms, but an OBE is not confined
to the original feasibility volume and can spuriously \emph{expand} the
solution set. Carlson instead works in a \emph{topological space}, where
property sets are static a~priori information that never expand, avoiding
the OBE pathology. Second, an \emph{application} move: STE is applied to
decryption for the first time.

\paragraph{The decryption reduction.} The solution space is a finite key
space $K$. Each intercepted ciphertext $c$ induces a property set of keys
under which $c$ decrypts to a meaningful message; the feasible key set is
the STE feasibility set of these property sets, and decryption succeeds
when it collapses to the true key --- the set theoretic form of Shannon's
unicity \citep{shannon1949}.

\paragraph{Named results (extracted from the dissertation text).}
\begin{itemize}
\item \textbf{Lemma 2.1:} the error $e$ (number of keys lacking an
  encryption property $\Phi_i$) is bounded, $0\le e\le|K|-1$.
\item \textbf{Lemma 2.2:} property sets behave as independent variables
  whose intersection is no larger than the smallest set; if
  $\Phi_j\subseteq\Phi_k$ then $\Phi_k$ is subsumed.
\item \textbf{Theorem 2.3:} each property set is an AEP typical set;
  hence information theory falls under the STE umbrella.
\item \textbf{Lemma 4.1:} a substitution ($S$) cipher on message $M$
  admits $(|A|-|T|)!$ isomorphic (equivalent) keys, where $T$ is the set
  of distinct symbols occurring in $M$.
\item \textbf{Theorem 4.1:} a permutation ($P$) cipher admits
  $\binom{|S_t|}{C}(|B|-|S_t|)!$ unique keys, $S_t$ the static-bit set.
\item \textbf{Theorem 5.2:} a permutation encryption reduces to a
  substitution encryption within a symbol's bit representation.
\item \textbf{Corollary 5.3:} a boundary-aligned PSP cipher is idempotent
  to a single $S$ cipher; therefore every block cipher is a block
  substitution cipher, so one attack (BCBB) suffices.
\end{itemize}

\paragraph{What we formalized.} The topological substrate ---
closedness and, via compactness and the finite intersection property,
existence of feasible points --- is \texttt{Ste.Topology}; it specializes
cleanly to Carlson's finite (discrete, hence compact) key spaces. The
STE-decryption skeleton (cipher systems, key property sets, feasible
keys, fairness of genuine traffic, unicity as ideal information,
observation monotonicity) is \texttt{Ste.Cipher}.

\section{Assessment: what is ripe to mechanize}

\begin{enumerate}
\item \textbf{Most tractable:} the cipher counting results (Lemma~4.1,
  Theorem~4.1). These are finite combinatorics over \texttt{Fin}-indexed
  symbol/bit alphabets; Mathlib's \texttt{Fintype.card\_perm}
  ($|\mathrm{Perm}\,\alpha| = (\#\alpha)!$) and factorial API supply the
  scaffolding.
\item \textbf{Structural:} the reduction chain (Theorem~5.2,
  Corollary~5.3) reducing all block ciphers to substitution ciphers ---
  an algebraic argument once ciphers are modeled as keyed permutations.
\item \textbf{Analytic:} POCS convergence (\citealp{youla1982,
  bauschke1996}) atop \texttt{Ste.Convex}, using Mathlib's Hilbert-space
  projection theory.
\item \textbf{Deepest:} Theorem~2.3 (property sets as AEP typical sets),
  which requires a probabilistic model and Mathlib's measure/information
  theory.
\end{enumerate}

\section{Distrust discipline}
Per project rule, no result here is asserted proven unless it compiles in
the \texttt{Ste} library checked by CI on every push to \texttt{main}.
The Carlson named results above are \emph{quoted from the source}, not
yet independently verified; they are targets, and their statuses live in
\texttt{log/2026-07-19-hypotheses-ste-core.md} (H9--H12, OPEN).

\bibliographystyle{plainnat}
\bibliography{../../refs}

\end{document}
